% !TEX root = ../../main.tex
\subsection{面向订单流程的多机制融合}

FoodShield 的机制创新主要体现在多种安全机制的组合使用。系统并未将匿名身份、订单认证、加密存储、哈希审计和条件溯源作为彼此独立的功能点，而是围绕外卖订单通信流程进行集成设计，使每个安全机制都服务于明确的业务环节。

首先，系统设计了基于 PID 的匿名身份机制。用户注册后由后端生成匿名身份标识 PID，PID 作为系统内部身份参与订单绑定和通信记录，而用户真实身份不直接暴露给骑手端。骑手端仅能看到订单相关信息和通信内容，无法通过普通业务接口获取用户真实身份或 PID，从而降低跨订单关联和身份暴露风险。

其次，系统设计了基于 HMAC-SM3 的订单 Token 认证机制。用户创建订单后，后端将订单号、匿名身份、时间戳等信息绑定生成 Token。用户进入订单通信房间前必须提交订单号和 Token，服务器验证通过后才允许其加入对应订单房间。该机制使用户能够在不暴露真实身份的情况下证明自己是合法订单参与方，避免仅凭订单号即可进入通信通道的越权风险。

再次，系统将 SM4 加密存储、SM3 消息摘要和 Merkle Tree 完整性验证结合起来。通信消息在写入数据库前进行加密处理，管理员端默认只查看消息哈希、订单摘要和 Merkle Root 等审计信息，而不是直接查看消息明文。系统以消息摘要作为 Merkle Tree 叶子节点，通过 Root 快照比对检测消息修改、删除、插入和重排等篡改行为，使聊天记录从普通数据库记录转变为可验证的审计对象。

最后，系统引入三端摘要一致性增强机制，将服务端、用户端和骑手端的 Merkle Root 进行比对。该机制提高了攻击者单独篡改平台侧日志和快照的难度，使日志完整性验证不再完全依赖单一服务端存储结果，从而进一步增强了通信记录的可信性。
